\subsection{Overall Solution Framework}
\label{subsec:framework}

The proposed method coordinates three decision blocks:
chamber assignment, sample-size search, and feasible scheduling.
For a sample vector \(\mathbf{s}=(s_p)_{p\in\mathcal{P}}\), each
chamber is assigned a temperature--humidity environment, producing
a configuration \(\mathbf{x}\).
Under this configuration, sample sizes are improved by local search.
Every candidate is decoded by an earliest-due-date (EDD)
constructive scheduler, which returns a feasible schedule \(\Pi\)
and the associated total tardiness \(T^{tot}\).
If local search does not improve the incumbent, a variable
neighborhood search (VNS) perturbs sample sizes and repeats the
same local search.
The inner procedures are specified in the following subsections;
this subsection states only the outer coordination.

A candidate is compared with the incumbent by the lexicographic
order
\begin{equation}
\label{eq:lex}
(T^{tot},\mathbf{1}^{\mathsf{T}}\mathbf{s})
\prec_{\mathrm{lex}}
(T^{tot}_{inc},\mathbf{1}^{\mathsf{T}}\mathbf{s}_{inc})
\end{equation}
if it has a strictly smaller total tardiness, or the same total
tardiness and a strictly smaller total sample count
\(\mathbf{1}^{\mathsf{T}}\mathbf{s}\).
The outer loop is required because a change in \(\mathbf{s}\)
alters the environment-wise workload---in particular pull-down
tests, which are executed once per sample---and may therefore
justify a different assignment \(\mathbf{x}\).
A new assignment, in turn, changes which products benefit from
additional samples.
Algorithm~\ref{alg:framework} states the resulting control logic.

\begin{algorithm}[H]
\caption{Overall Solution Framework}
\label{alg:framework}
\begin{algorithmic}[1]
\State \textbf{Input:} instance data
\State \textbf{Output:} \((\mathbf{x}_{inc},\mathbf{s}_{inc},\Pi_{inc},T^{tot}_{inc})\)
\State Set \(s_p\leftarrow s_{\min}\) for all \(p\in\mathcal{P}\); \(\mathbf{s}\leftarrow(s_p)_{p\in\mathcal{P}}\) \label{step:fw-init}
\State \((\mathbf{x}_{inc},\mathbf{s}_{inc},T^{tot}_{inc})\leftarrow(\emptyset,\mathbf{s},+\infty)\)
\For{\(i\leftarrow 1\) to \(I_{\max}\)}
    \State \(\mathbf{x}\leftarrow\) \Call{ChamberAssignment}{\(\mathbf{s}\)} \label{step:fw-assign}
    \State \((\mathbf{s}_{cand},\Pi_{cand},T^{tot}_{cand})\leftarrow\) \Call{SampleLocalSearch}{\(\mathbf{x},\mathbf{s}\)} \label{step:fw-ls}
    \If{\((T^{tot}_{cand},\mathbf{1}^{\mathsf{T}}\mathbf{s}_{cand})\prec_{\mathrm{lex}}(T^{tot}_{inc},\mathbf{1}^{\mathsf{T}}\mathbf{s}_{inc})\)} \label{step:fw-lex}
        \State \((\mathbf{x}_{inc},\mathbf{s}_{inc},\Pi_{inc},T^{tot}_{inc})\leftarrow(\mathbf{x},\mathbf{s}_{cand},\Pi_{cand},T^{tot}_{cand})\)
        \State \(\mathbf{s}\leftarrow\mathbf{s}_{cand}\)
    \Else
        \State \((\mathbf{s}_{cand},\Pi_{cand},T^{tot}_{cand})\leftarrow\) \Call{VNS}{\(\mathbf{x},\mathbf{s}_{inc},\Pi_{inc},T^{tot}_{inc}\)} \label{step:fw-vns}
        \If{\((T^{tot}_{cand},\mathbf{1}^{\mathsf{T}}\mathbf{s}_{cand})\prec_{\mathrm{lex}}(T^{tot}_{inc},\mathbf{1}^{\mathsf{T}}\mathbf{s}_{inc})\)}
            \State \((\mathbf{x}_{inc},\mathbf{s}_{inc},\Pi_{inc},T^{tot}_{inc})\leftarrow(\mathbf{x},\mathbf{s}_{cand},\Pi_{cand},T^{tot}_{cand})\)
            \State \(\mathbf{s}\leftarrow\mathbf{s}_{cand}\)
        \Else
            \State \Return \((\mathbf{x}_{inc},\mathbf{s}_{inc},\Pi_{inc},T^{tot}_{inc})\) \label{step:fw-stop}
        \EndIf
    \EndIf
\EndFor
\State \Return \((\mathbf{x}_{inc},\mathbf{s}_{inc},\Pi_{inc},T^{tot}_{inc})\)
\end{algorithmic}
\end{algorithm}

The search starts from the minimum sample policy
(Step~\ref{step:fw-init}).
Each outer iteration assigns chambers to the current \(\mathbf{s}\)
(Step~\ref{step:fw-assign}) and then applies sample-size local
search under that assignment (Step~\ref{step:fw-ls}).
If the resulting candidate is lexicographically better than the
incumbent (Step~\ref{step:fw-lex}), it is accepted and the next
iteration reassigns chambers to the updated sample vector.
Otherwise VNS is called from the incumbent
(Step~\ref{step:fw-vns}).
The procedure stops when VNS fails to improve the incumbent
(Step~\ref{step:fw-stop}) or when the iteration limit
\(I_{\max}\) is reached.
